% VERSION MAP: v0.5 method-first (09-01) | v0.6 question-first + genealogy (09-03 Fable) | v0.61 narrative-first (09-03 Sol) | v0.7 = Sol spine + Fable body, PI rulings (09-03) | v0.71 = Sol red-team R2 adjudicated (09-03) | v0.8 = SOTA comparison filled (09-05) | v0.81 = PathSelect row = hier, flat as ablation, own A1/A3 rows under comparison protocol, paired 10-fold statistics; OPCM sentence corrected (09-05) | v0.82 = Sol R3 adjudicated: quoted-vs-rerun framing, no parity/SOTA wording, memory storage boundary, controlled-variants block, 2x2 novelty framing (09-06) | v0.83 = PI restructure: main table = 12 quoted + PathSelect only; ablations rebuilt as leave-out narrative; six-question and lifecycle tables removed; codes stripped (09-06) | v0.84 = forward-order results and order-dependent hierarchy reported (09-06) | v0.85 = forward-order leave-out chain; robustness synthesis (09-06) | v0.86 = QPMIL-VL reproduction check with released code, both orders (09-07); v0.87 | v0.88 = leave-out and order tables merged; float barrier before Discussion (09-07) | v0.89 = Fig.1 v2.0 integrated (09-08); v0.89a = batch-size deviation disclosed; quoting-precedent sentence made precise (09-08)
\documentclass{article}
\usepackage{spconf,amsmath,amssymb,graphicx,hyperref}
\usepackage{booktabs}
\usepackage{multirow}
\usepackage{placeins}

\def\x{{\mathbf x}}
\def\L{{\cal L}}
\newcommand{\Fg}{F_g}
\newcommand{\Fp}{F_p}
\newcommand{\Psel}{\mathcal{P}}
\newcommand{\Mem}{\mathcal{M}}
\newcommand{\pending}[1]{\textbf{[#1]}}

\title{REMEMBERING WHERE TO LOOK: CONTINUAL HIERARCHICAL EVIDENCE SELECTION FOR WHOLE-SLIDE IMAGES}

\name{Hsi-Ren Hung \qquad [Second Author] \qquad Huei-Fang Yang\sthanks{corresponding author. This work was supported in part by the National Science and Technology Council, Taiwan under Grant NSTC 112-2221-E-110-047-MY3, 114-2640-E-110-004, and 114-2634-F-006-002.}}
\address{Department of Computer Science and Engineering, National Sun Yat-sen University, Taiwan}

\begin{document}
\ninept
\maketitle

% -----------------------------------------------------------------------------
\begin{abstract}
Whole-slide image (WSI) diagnosis is an evidence-acquisition problem: a model must select a few patches from thousands, and its prediction is only as reliable as the evidence it chooses.
Continual-learning (CL) methods usually preserve classifiers or representations, but a budgeted diagnostic system can fail in a different way: its diagnostic pathway may remain intact while its selector gradually forgets \emph{where to look}.
We introduce \textbf{PathSelect}, which makes the evidence-selection policy itself the object of continual learning.
To isolate this failure, the image encoder, text encoder, and cosine diagnostic head are frozen; only a hierarchical group-to-patch selector is learned.
Within a task, three supervision signals answer complementary questions: whether the selected evidence predicts correctly, whether it is semantically discriminative, and which candidate patch would counterfactually reduce diagnostic loss.
Across tasks, per-task LoRA residuals are merged exactly into one shared selector, while a feature-free Selection Memory preserves past behavior through dual-level distillation, a utility-preservation hinge, and replay.
The key distinction is between \emph{behavioral continuity}---how the selector scores evidence---and \emph{functional continuity}---whether newly selected evidence remains diagnostically useful.
On four sequential TCGA tasks, PathSelect raises final class-incremental accuracy by 34.6 points over sequential fine-tuning in every seed, halves forgetting relative to replay alone ($0.0539$ vs.\ $0.1102$), and retains two orders of magnitude more of its originally selected evidence; under QPMIL-VL's reverse-order splits and evaluation protocol it reaches 0.854 ACC, within 0.5 points of QPMIL-VL's reported 0.859, while adapting 16.4k selector parameters per task (QPMIL-VL reports 0.365M) and storing no features, with inference on one task-free shared selector.
\end{abstract}

\begin{keywords}
whole-slide images, continual learning, evidence selection, computational pathology, LoRA
\end{keywords}

% -----------------------------------------------------------------------------
\section{Introduction}
\label{sec:intro}

A gigapixel WSI can contain tens of thousands of patches, yet a diagnosis often depends on a small subset.
Practical systems therefore operate under an \emph{evidence budget}: choose $B$ patches, then decide.
This creates a learning problem that is easy to miss in conventional continual learning.
A model may preserve its classifier and representation reasonably well, yet after learning new disease cohorts it can still fail because it no longer chooses the tissue that once supported the correct decision.

We therefore ask a narrower question:
\begin{quote}
\emph{When diagnostic tasks arrive sequentially, does a model forget where to look, and can that ability be preserved, and its loss measured, directly?}
\end{quote}
PathSelect is designed around this question rather than around a collection of modules.
Its first principle is \textbf{attribution by freezing}.
CONCH image/text encoders \cite{lu2024conch} and the cosine diagnostic rule are fixed throughout training.
There is no trainable classifier.
The only learned object is an evidence-selection policy.
For paired methods evaluated at the same checkpoint and class set the diagnostic pathway is identical, so between-method differences arise from selector updates; within-method class-incremental changes over time additionally include label-space expansion. Direct selection metrics provide an additional head-independent measurement.

Our hierarchical substrate follows the coarse-to-fine tissue-aware idea of HistoSelect \cite{huang2026histoselect}: first allocate attention among tissue groups, then select patches within each group.
We do \emph{not} claim this hierarchy itself as our primary novelty.
Instead, PathSelect starts where such selectors normally stop: after a selector has learned where to look, how can one keep that ability when tasks arrive sequentially?

This perspective leads to a deliberate separation between two kinds of learning.
Within a task, PathSelect must \emph{learn where to look now}.
Across tasks, it must \emph{keep looking well}.
The first is trained by diagnostic, semantic, and counterfactual-utility supervision.
The second is trained by three preservation questions: (i) do I still score old evidence similarly, (ii) even if I select differently, is the new evidence at least as useful, and (iii) can I still solve the old diagnostic problem?

\textbf{Contributions.}
\begin{itemize}
    \item We formulate continual WSI diagnosis with the \emph{evidence-selection policy}---rather than a classifier---as the learned, preserved, and measured object, isolated by a frozen diagnostic pathway; to our knowledge, prior WSI continual-learning methods adapt MIL, attention, or localization components, whereas recent budgeted evidence-selection and active-perception methods are not continual; PathSelect isolates the intersection---the diagnostic pathway is fixed, and a hard, budgeted evidence selector is the only learned and directly evaluated continual object.
    \item We introduce a fixed-cost continual protocol in which per-task LoRA residuals are merged into one shared hierarchical selector, while a feature-free Selection Memory preserves both \emph{selection behavior} and \emph{selection usefulness}; under QPMIL-VL's splits and metrics it reaches 0.854 ACC, within 0.5 points of the strongest published result in that comparison (0.859), while adapting 16.4k selector parameters per task against QPMIL-VL's reported 0.365M ($22.3\times$ fewer; different adapted components, reported as adaptation-cost context).
    \item We pair the preservation objectives with selection-level measurements: Jaccard for identity drift and utility retention for functional drift, and show experimentally that these two notions are related but not interchangeable.
\end{itemize}

% -----------------------------------------------------------------------------
\section{Related Work}
\label{sec:related}

\textbf{Evidence selection in WSIs.}
MIL methods such as ABMIL, CLAM, and TransMIL \cite{ilse2018abmil,lu2021clam,shao2021transmil} aggregate patch features with learned attention, but selection is entangled with a trained diagnostic model.
HistoSelect \cite{huang2026histoselect} makes selection explicitly hierarchical, using tissue-aware coarse-to-fine reasoning.
PathSelect inherits this group$\rightarrow$patch inductive bias but changes the scientific setting: the diagnostic pathway is frozen and the selector is learned under a sequential task stream. HistoSelect conditions selection on a per-question query for VLM reasoning; PathSelect instead fixes a task-free tissue vocabulary, because a policy that must persist across tasks needs a grouping that does not move.

\textbf{LoRA-based continual learning.}
Low-rank adaptation \cite{hu2022lora} provides parameter-efficient task updates.
Prior CL methods use low-rank or modular updates to preserve classifier/backbone knowledge or retrieve task-specific components \cite{liang2024inflora,wu2025sdlora,he2025cllora,chen2026flexem}.
PathSelect instead applies continual learning only to the evidence selector.
Importantly, LoRA merging is used as a \emph{deployment substrate}, not claimed as the preservation mechanism: our experiments show that merging alone does not prevent forgetting.

\textbf{Continual WSI learning.} ConSlide \cite{huang2023conslide} introduced rehearsal-based continual WSI analysis on a four-organ TCGA sequence; MICIL \cite{meseguer2024micil} combined MIL with class-incremental learning; QPMIL-VL \cite{gou2025qpmil} brought vision--language prototypes to incremental WSI classification on CONCH features; and, closest to us, Li et al.\ \cite{li2025akd} located forgetting in the attention layers of MIL models and proposed attention distillation with a pseudo-bag memory of informative patches. CoMEL \cite{lee2025comel} adapts MIL aggregators continually with orthogonal weighted low-rank adaptation, and MergeSlide \cite{bui2026mergeslide} merges per-task fine-tuned models by orthogonal projection with task-to-class prompt inference. All of these retain a trainable MIL diagnostic model in which attention or aggregation is part of the learned diagnostic function. PathSelect differs on four axes: the entire diagnostic pathway is frozen and the selector is the sole learned object; selection is a hard, budgeted set rather than soft attention weights; preservation targets the utility of the selected set in addition to score behavior, so identity may change without functional regression; and selection identity and utility are measured directly rather than inferred from accuracy. Conversely, budgeted evidence selection \cite{huang2026histoselect} and active-perception navigation over WSIs \cite{chen2026adaptivepath} learn where to look but are not continual; PathSelect sits at the intersection of the two lines.

% -----------------------------------------------------------------------------
\section{Method}
\label{sec:method}

\subsection{Study object: a frozen diagnostic pathway}
\label{ssec:setup}

Tasks $\tau=1,\ldots,T$ arrive sequentially.
For a slide, frozen CONCH features are
\[
\{x_i\}_{i=1}^{N}, \qquad x_i\in\mathbb{R}^{512}.
\]
The selector must choose an evidence set
\[
\Psel=\{x_{i_1},\ldots,x_{i_B}\}, \qquad |\Psel|=B=8.
\]
Only this evidence reaches the diagnostic rule; the budget $B{=}8$ is the operating point at which learned selection peaked over random selection in a preliminary sweep (Sec.~\ref{sec:exp}).
If selector scores for the chosen patches are $s_i$, the WSI representation is the $\ell_2$-normalized score-weighted mean $z=\sum_{i\in\Psel}\mathrm{softmax}(s)_i\,x_i/\|\cdot\|_2$.
Prediction is the cosine similarity between $z$ and frozen class-text embeddings $\{t_c\}$, followed by softmax.
Training is task-aware; inference is task-free over all seen classes.
No task identity, adapter bank, memory lookup, or trained diagnostic head is used at test time.

\begin{figure*}[t]
    \centering
    \IfFileExists{figures/fig1_architecture_v2_0.pdf}{
        \includegraphics[width=\textwidth]{figures/fig1_architecture_v2_0.pdf}
    }{
        \fbox{\parbox[c][2.0cm][c]{0.96\textwidth}{\centering
        \textbf{Architecture figure placeholder.}\\
        Frozen CONCH features $\rightarrow$ semantic tissue groups $\rightarrow$
        group selector $\Fg$ $\rightarrow$ whole-budget quotas $\rightarrow$
        patch selector $\Fp$ $\rightarrow$ $B{=}8$ evidence $\rightarrow$
        frozen cosine diagnosis;\\
        across tasks: task LoRA $\rightarrow$ exact merge into shared selector;
        Selection Memory $\rightarrow$ KD / utility-equivalence / replay.}}
    }
    \caption{\textbf{PathSelect overview.} Top (solid arrows; training and inference): frozen CONCH patch features are grouped by eight fixed tissue prompts; the group selector $\Fg$ turns group scores into integer quotas ($\sum_j b_j=B=8$), the patch selector $\Fp$ keeps the top $b_j$ patches within each group, and the frozen head diagnoses from the selected evidence alone by cosine with the class-text embeddings---inference ends there. Bottom (dashed; training only): each task trains fresh LoRA adapters on the same two MLPs and folds them into one shared selector at the task end; a Selection Memory of 512 behavioral snapshots (scores, indices, utilities; no images or features) drives replay, dual-level distillation, and the utility-preservation hinge.}
    \label{fig:arch}
\end{figure*}

\subsection{Hierarchical budgeted selection}
\label{ssec:hier}

Following the tissue-aware coarse-to-fine principle of HistoSelect \cite{huang2026histoselect}, patches are assigned semantically rather than clustered.
Eight fixed tissue prompts are encoded once with the frozen text encoder.
Each patch is assigned by
\[
G_j=\left\{x_i:
j=\arg\max_k \cos(x_i,p_k)\right\},
\]
and each group prototype is
\[
g_j=\operatorname{mean}(x_i\in G_j).
\]

Two small two-layer MLPs ($\Fg$ and $\Fp$, identical in structure and independent in weights) implement the policy; their input is zero-padded to a common width so that every ablation arm shares one architecture and parameter count.
The group selector $\Fg$ produces group scores $r_j$.
A largest-remainder allocation converts $\mathrm{softmax}(r)$ into integer quotas
\[
b_j\in\mathbb{N}, \qquad \sum_j b_j=B.
\]
The patch selector $\Fp$ produces patch scores $s_i$, and the top $b_j$ patches are selected within each group.
Thus the hierarchy answers two different questions:
\emph{which tissue groups deserve the budget?} and \emph{which patches within those groups deserve the slots?}

\subsection{Learning where to look: three within-task supervision signals}
\label{ssec:within}

The within-task objective is
\begin{equation}
\L_{\mathrm{ev}}
=
\L_{\mathrm{diag}}
+\beta_s\L_{\mathrm{sem}}
+\beta_u\L_{\mathrm{util}}.
\end{equation}
These terms are not redundant; they provide supervision at different levels.

\textbf{Outcome supervision: $\L_{\mathrm{diag}}$.}
Cross-entropy on the frozen diagnostic logits asks:
\emph{did the final evidence set answer correctly?}
It supervises the joint outcome but provides weak credit assignment among many candidate patches.

\textbf{Semantic supervision: $\L_{\mathrm{sem}}$.}
A frozen class-text prior scores how diagnostically discriminative a patch appears in the CONCH embedding space. (At inference the selector receives no class text, task identity, or label; during training, class-derived information enters only as supervision targets---this prior and $\L_{\rm diag}$. The attribution claim concerns the frozen non-selector parameters, not the absence of class-derived supervision.)
The selector distribution is softly aligned to this prior with KL divergence.
This is a weak semantic anchor, not the final diagnostic target.

\textbf{Counterfactual credit: $\L_{\mathrm{util}}$.}
For candidate patch $i$, define a marginal counterfactual gain
\begin{equation}
u_i
=
\L_{\mathrm{diag}}(E)
-
\L_{\mathrm{diag}}(E\cup\{x_i\}),
\label{eq:ui}
\end{equation}
where $E$ denotes the current evidence context.
Positive $u_i$ means that adding the candidate would reduce diagnostic loss; negative $u_i$ means it would be harmful.
The gains are treated as a no-gradient teacher distribution, and $\L_{\mathrm{util}}=\mathrm{KL}(\mathrm{softmax}(u)\,\Vert\,\mathrm{softmax}(s))$ trains the patch selector to rank candidates by measured usefulness.
Lowercase $u_i$ is therefore a \emph{candidate-level} teaching signal.

The value of the \emph{whole selected evidence set} is a different quantity:
\begin{equation}
U(\Psel)
=
\log C-\mathrm{CE}(\Psel,y).
\label{eq:bigU}
\end{equation}
It asks how much the complete selected set improves over an empty/uniform diagnostic state.
This distinction---candidate utility $u_i$ versus set utility $U$---becomes central in continual learning.


\subsection{One shared selector: LoRA consolidation}
\label{ssec:lora}

For task $\tau$, low-rank LoRA adapters are attached to the four linear layers of $\Fg$ and $\Fp$.
The base selector is frozen during the task and only the current adapters train.
At the task boundary, $W_{\tau}=W_{\tau-1}+B_\tau A_\tau$, which is an exact rewrite of the effective weight used immediately before merging.
The adapter is then reset for the next task.
Thus parameter count does not grow with $T$ and inference always uses one shared selector.

Crucially, \textbf{merging is not the anti-forgetting mechanism}.
The merge-only arm performs no better systematically than naive sequential fine-tuning.
LoRA supplies a clean, fixed-cost way to accumulate task updates; preservation requires additional constraints.

\subsection{Keeping looking well: preserve behavior and function}
\label{ssec:cl}

After a task is learned, PathSelect stores a bounded Selection Memory
$\Mem$ of at most 512 behavioral snapshots using reservoir sampling.
A snapshot stores an internal sample reference, the task tag, group scores $r_{\rm old}$ (8 values), candidate indices ($\le$256), patch scores $s_{\rm old}$ ($\le$256), and utilities ($\le$256); it contains no raw image or feature tensor. The 512-entry memory occupies 3.11\,MB when serialized. Replay nevertheless assumes authorized training-time access to the corresponding pre-extracted feature records, whose storage ($19.5$\,GB for this benchmark) is not included in that figure.

For each current sample, one old snapshot is replayed.
The continual objective is
\begin{align}
\L_{\mathrm{CL}}
={}&
\lambda_{\mathrm{kd}}
\left[
\mathrm{KL}(r_{\mathrm{old}}\Vert r_{\mathrm{new}})
+
\mathrm{KL}(s_{\mathrm{old}}\Vert s_{\mathrm{new}})
\right]
\nonumber\\
&+
\lambda_{\mathrm{eq}}
\max(0,U_{\mathrm{old}}-U_{\mathrm{new}})
+
\lambda_r \L_{\mathrm{diag}}^{\mathrm{rep}},
\label{eq:lcl}
\end{align}
with all three $\lambda$'s set to 1.

Equation~\ref{eq:lcl} deliberately separates three preservation questions.

\textbf{Behavioral continuity: $\L_{\rm KD}$.}
Dual-level distillation preserves the old \emph{soft selection policy}: group preferences and patch-score distributions.
It does not require the final hard top-$B$ set to be identical.

\textbf{Functional continuity: $\L_{\rm eq}$.}
Selection identity should be allowed to change if an alternative evidence set is equally or more useful.
The hinge
\[
\L_{\rm eq}=\max(0,U_{\rm old}-U_{\rm new})
\]
therefore acts as a \emph{floor, not a ceiling}: regression is penalized, improvement is free.
This is the central difference between ``remember the same evidence'' and ``retain the same evidence-selection competence.''

\textbf{Task competence: replay.}
$\L_{\mathrm{diag}}^{\mathrm{rep}}$ is simply the ordinary diagnostic loss on a replayed old slide.
Replay is a data mechanism rather than a new mathematical objective.
All within-task and continual terms update the \emph{same current selector LoRA} in one backward pass:
\[
\L_{\mathrm{total}}=\L_{\mathrm{ev}}+\L_{\mathrm{CL}}.
\]

% -----------------------------------------------------------------------------
\section{Experiments}
\label{sec:exp}

\subsection{Protocol and measurements}

We use the four-cohort TCGA benchmark introduced by ConSlide \cite{huang2023conslide} in the form released by QPMIL-VL \cite{gou2025qpmil}: CLAM-tiled $256\times256$ patches at $10\times$, pre-extracted CONCH features, and ten patient-level folds. Tasks arrive in the reverse order ESCA$\rightarrow$RCC$\rightarrow$BRCA$\rightarrow$LUNG, each a binary subtype task, for a final eight-class label space. Two protocols are used. The \emph{comparison protocol} follows QPMIL-VL: ten folds, one run per fold with the seed equal to the fold index, and the metrics average accuracy (ACC), masked ACC (task-IL), forgetting \cite{chaudhry2018rwalk}, and backward transfer (BWT) \cite{lopezpaz2017gem}. The \emph{ablation protocol} fixes fold 1 (2,273 training slides; test sizes 15/76/93/95) and varies five seeds with paired win counts.
All runs use 5 epochs per task, learning rate $10^{-3}$, $B=8$, LoRA rank 4, $\beta_s{=}\beta_u{=}0.1$, unit $\lambda$'s, and seeds 0--4. Every comparison was registered with a two-sided decision rule before the corresponding runs: win counts over the five seeds decide, the incumbent design is retained on a split, and one registered expectation that the data refuted is reported as measured (Sec.~\ref{ssec:ablation}).

The controlled variants are sequential full fine-tuning of the selector (SeqFT), LoRA with boundary merging only, replay, replay+distillation, and full PathSelect (replay+distillation+utility hinge).
Per-task specialists and joint offline training are references rather than deployable continual baselines.

We measure outcome and selection separately.
Class-IL accuracy evaluates over all seen classes; task-IL accuracy restricts prediction to the task's own pair; leakage is the fraction predicted into another task's class set.
Selection Jaccard, $|\Psel_{\rm own}\cap\Psel_{\rm final}|/|\Psel_{\rm own}\cup\Psel_{\rm final}|$, measures hard identity retention between the set selected immediately after learning a task and the set selected at the end.
Jaccard is \emph{evaluation only}; it is not the distillation loss.
Functional retention is measured by
\[
\Delta U=U_{\rm final}-U_{\rm own}.
\]
With five seeds, comparisons are paired by seed and accompanied by win counts; 5/5 is called systematic, 4/5 directional, and no $p$-values are reported at this seed count.

\subsection{Comparison with continual WSI methods}
\label{ssec:sota}

\begin{table*}[t]
\centering\footnotesize
\caption{\textbf{Comparison under QPMIL-VL's reverse-order splits and evaluation protocol} (ESCA$\rightarrow$RCC$\rightarrow$BRCA$\rightarrow$NSCLC, ten-fold CV, mean$\pm$sd). Published rows are quoted from QPMIL-VL's Table~2 \cite{gou2025qpmil} and were not rerun by us; PathSelect uses the identical public split files and metric definitions, while training implementations and hyperparameters follow each method and are therefore not controlled across papers. Unlike the cited MIL baselines, PathSelect adapts only the evidence selector and keeps the diagnostic pathway frozen, so parameter counts characterize adaptation cost rather than capacity-matched architectures; external rows provide published benchmark context rather than a controlled reimplementation comparison. Buffer sizes in WSIs; PathSelect's Selection Memory holds 512 behavioral snapshots (3.1\,MB, excluding the feature store required for replay). Forward-order results are in Table~\ref{tab:chain}. $^\ddagger$In the forward order, which the released code runs natively, our rerun reproduces the published result to within $0.6$ points ($0.884{\pm}.033$ vs.\ $0.890{\pm}.021$; forgetting $0.038$ vs.\ $0.027$), which cross-validates the data, folds, and metric definitions. The released code fixes the task order in a class-ensemble file, so the reverse order required reordering that file, and our reverse rerun used the code's default mini-batch of 16 where the paper reports 8 (both settings-level deviations); the published reverse value lies within one standard deviation of our reproduction, and we report both rather than replace one with the other. Mixing rerun and directly quoted published entries is also the practice of QPMIL-VL's own tables, which mark quoted entries in gray.}
\label{tab:sota}
\begin{tabular}{llccccc}
\toprule
Type & Method & ACC$\uparrow$ & Forgetting$\downarrow$ & BWT$\uparrow$ & Masked ACC$\uparrow$ & Buffer \\
\midrule
\multicolumn{7}{l}{\emph{Published methods---values quoted from \cite{gou2025qpmil}}} \\
Reference & JointTrain (offline reference) & $0.908{\pm}.022$ & --- & --- & $0.937{\pm}.022$ & --- \\
Reference & FineTune (lower) & $0.234{\pm}.008$ & $0.927{\pm}.023$ & $-0.927{\pm}.023$ & $0.803{\pm}.060$ & 0 \\
Regularization & EWC \cite{kirkpatrick2017ewc} & $0.235{\pm}.011$ & $0.928{\pm}.020$ & $-0.928{\pm}.020$ & $0.833{\pm}.069$ & 0 \\
Regularization & LwF \cite{li2017lwf} & $0.236{\pm}.016$ & $0.908{\pm}.030$ & $-0.908{\pm}.030$ & $0.900{\pm}.041$ & 0 \\
Rehearsal & A-GEM \cite{chaudhry2019agem} & $0.536{\pm}.047$ & $0.527{\pm}.067$ & $-0.527{\pm}.067$ & $0.872{\pm}.026$ & 30 \\
Rehearsal & ER-ACE \cite{caccia2022erace} & $0.703{\pm}.049$ & $0.281{\pm}.062$ & $-0.279{\pm}.063$ & $0.889{\pm}.041$ & 30 \\
Rehearsal & DER++ \cite{buzzega2020derpp} & $0.684{\pm}.055$ & $0.310{\pm}.069$ & $-0.307{\pm}.072$ & $0.910{\pm}.043$ & 30 \\
Rehearsal & ER \cite{chaudhry2019er} & $0.644{\pm}.028$ & $0.387{\pm}.050$ & $-0.386{\pm}.049$ & $0.901{\pm}.035$ & 30 \\
WSI CL & ConSlide \cite{huang2023conslide} & $0.499{\pm}.025$ & $0.058{\pm}.032$ & $-0.021{\pm}.039$ & $0.854{\pm}.039$ & 30 \\
VLM & AttriCLIP \cite{wang2023attriclip} & $0.694{\pm}.058$ & $0.207{\pm}.063$ & $-0.207{\pm}.063$ & $0.861{\pm}.019$ & 0 \\
VLM & MI-Zero \cite{lu2023mizero} & $0.839{\pm}.034$ & --- & --- & $0.909{\pm}.019$ & 0 \\
VLM WSI CL & QPMIL-VL \cite{gou2025qpmil} & $0.859{\pm}.032$ & $0.064{\pm}.031$ & $-0.064{\pm}.031$ & $0.925{\pm}.018$ & 0 \\
\midrule
\multicolumn{7}{l}{\emph{Reproduction check---our run of the released QPMIL-VL code on the same folds (one RTX 4090)$^\ddagger$}} \\
VLM WSI CL & QPMIL-VL (reproduced) & $0.814{\pm}.065$ & $0.147{\pm}.091$ & $-0.147{\pm}.091$ & $0.908{\pm}.043$ & 0 \\
\midrule
\multicolumn{7}{l}{\emph{Our method---run under the same splits and metrics}} \\
Selection CL & PathSelect (ours) & $0.854{\pm}.044$ & $0.060{\pm}.041$ & $-0.056{\pm}.038$ & $0.925{\pm}.019$ & 512 snap.\ (3.1\,MB) \\
\bottomrule
\end{tabular}
\end{table*}

Table~\ref{tab:sota} places PathSelect in the continual WSI landscape under QPMIL-VL's splits and evaluation protocol. PathSelect reaches $0.854{\pm}.044$ ACC, $0.925{\pm}.019$ masked ACC, and $0.060{\pm}.041$ forgetting: the second-highest ACC among the continual methods in the table, $0.005$ below QPMIL-VL's reported $0.859$, and above every regularization, rehearsal, ConSlide, and prompt-based row; the offline joint-training reference reaches $0.908$. Rerunning QPMIL-VL with its released code reproduces its forward-order result to within $0.6$ points ($0.884{\pm}.033$ vs.\ $0.890{\pm}.021$) but gives $0.814{\pm}.065$ ACC and $0.147{\pm}.091$ forgetting in the reverse order (fold range $0.729$--$0.900$), where the code required a settings-level change; relative to this reproduction PathSelect is $4.0$ points higher and forgets less than half as much, and relative to the published value it is $0.5$ points lower. We therefore describe PathSelect as competitive with the strongest published method rather than above or below it. ConSlide reports slightly lower forgetting ($0.058$ vs.\ $0.060$) at substantially lower ACC ($0.499$), so we do not claim the lowest forgetting. PathSelect adapts 16.4k selector parameters per task, versus QPMIL-VL's reported $0.365$M learnable parameters ($22.3\times$ fewer) and $0.5$--$39$M for the MIL-based rows; these counts refer to different adapted components and are reported as adaptation-cost context, together with the absence of a prototype pool or task-specific module at inference and a 3.1\,MB behavioral memory (a 30-WSI rehearsal buffer of features would occupy about $180$\,MB at $6.1$\,MB per slide).

\subsection{What each part of PathSelect contributes}
\label{ssec:ablation}


Learned selection is itself meaningful: in a preliminary budget sweep with a learned flat selector (inference only; seven budgets $\times$ four tasks), learned selection beat random selection in all 28 cells and peaked at $B=8$ ($0.8797$ vs.\ $0.6607$).
Table~\ref{tab:chain} removes one part of PathSelect at a time under the comparison protocol, in both task orders.
The largest loss comes from removing preservation altogether: fine-tuning only the selector drops ACC from $0.830$ to $0.452$ and raises forgetting from $0.090$ to $0.601$, in every fold on every metric---while still beating full-model fine-tuning under the same splits ($0.234$ in Table~\ref{tab:sota}), so the frozen pathway alone halves the collapse.
Replay recovers most of the accuracy ($0.821$); distillation and the utility hinge then add $0.8$ ACC points (7/10 folds) and remove $2.4$ forgetting points (8/10); the hierarchy adds $2.5$ ACC points (8/10) and removes $2.9$ forgetting points (7/10) under the canonical reverse order. Under the forward order (QPMIL-VL's Table~1 protocol; Table~\ref{tab:chain}) the chain is $0.529\rightarrow0.841$ (replay) $\rightarrow0.841$ (full, flat) $\rightarrow0.811$ (full, hierarchical): preservation is again decisive (10/10 folds on every metric), distillation and the hinge no longer add accuracy (5/10) but still reduce forgetting ($-0.9$ points, 8/10), and the hierarchy costs $3.0$ points (9/10). Across both orders, what is robust is the preservation framework itself and the forgetting reduction of the two behavioral terms; their accuracy increment and the substrate's benefit are order-dependent, and we report both rather than select an order. The hierarchy is retained as the pre-registered choice on the canonical order.
Without any learned selection the frozen head reaches $0.679$, so the policy contributes $17.5$ points.

\begin{table}[t]
\centering\footnotesize
\setlength{\tabcolsep}{4pt}
\caption{\textbf{Leave-out chain in both task orders} (comparison protocol, ten-fold CV, mean$\pm$sd). Reverse = ESCA$\rightarrow$RCC$\rightarrow$BRCA$\rightarrow$LUNG; forward = the opposite. Fold-paired win counts are given in the text. QPMIL-VL: values reported in \cite{gou2025qpmil} and our rerun of the released code.}
\label{tab:chain}
\resizebox{\columnwidth}{!}{%
\begin{tabular}{lcccc}
\toprule
 & \multicolumn{2}{c}{Reverse} & \multicolumn{2}{c}{Forward} \\
\cmidrule(lr){2-3}\cmidrule(lr){4-5}
Configuration & ACC$\uparrow$ & Forg.$\downarrow$ & ACC$\uparrow$ & Forg.$\downarrow$ \\
\midrule
Selector fine-tuning (no preservation) & $.452{\pm}.071$ & $.601$ & $.529{\pm}.077$ & $.503$ \\
\ + replay of snapshots & $.821{\pm}.036$ & $.114$ & $.841{\pm}.019$ & $.064$ \\
\ + distillation + utility hinge (flat) & $.830{\pm}.029$ & $.090$ & $.841{\pm}.043$ & $.055$ \\
\ + hierarchical selector (PathSelect) & $.854{\pm}.044$ & $.060$ & $.811{\pm}.029$ & $.061$ \\
Frozen head, mean pooling (no selection) & $.679{\pm}.032$ & --- & --- & --- \\
\midrule
QPMIL-VL (reported) & $.859{\pm}.032$ & $.064$ & $.890{\pm}.021$ & $.027$ \\
QPMIL-VL (reproduced, released code) & $.814{\pm}.065$ & $.147$ & $.884{\pm}.033$ & $.038$ \\
\bottomrule
\end{tabular}}
\end{table}

\begin{table}[t]
\centering
\footnotesize
\caption{\textbf{Preservation-component ablation} (flat, fold 1, 5-seed means). Replay restores task attribution, the utility hinge assigns credit, distillation stabilizes scoring; no single term or pair matches the full set.}
\label{tab:ablation}
\resizebox{0.95\columnwidth}{!}{%
\begin{tabular}{ccc|ccc}
\toprule
Replay & Distill. & Utility & Class-IL$\uparrow$ & Task-IL$\uparrow$ & Leak.$\downarrow$\\
\midrule
 & & & .447 & .811 & .476\\
 & \checkmark & & .600 & .853 & .323\\
 & & \checkmark & .809 & .885 & .091\\
\checkmark & & & .778 & .907 & .104\\
\checkmark & \checkmark & & .797 & .897 & .117\\
\checkmark & \checkmark & \checkmark & \textbf{.824} & \textbf{.915} & \textbf{.101}\\
\bottomrule
\end{tabular}}
\end{table}

Table~\ref{tab:ablation} separates the three preservation losses on the ablation protocol.
Replay is the largest single contributor ($0.447\rightarrow0.778$) because it restores task attribution: with distillation alone, leakage stays at $0.323$, three times the full method---in one seed 59 of 76 renal slides land in the lung classes---so distillation preserves \emph{how} the selector scores but not \emph{which task} the evidence belongs to.
The utility hinge alone reaches $0.809$; its increment on top of replay and distillation is order-sensitive and we report it as such.
Under the hierarchy, removing the group level of distillation costs $3.95$ task-IL and $3.71$ class-IL points (5/5), so quota-level and patch-level preservation act at different resolutions.

\begin{table}[t]
\centering\footnotesize
\caption{\textbf{Selection-level metrics under the ablation protocol} (flat, fold 1, mean over 5 seeds): what accuracy alone cannot show.}
\label{tab:selmetrics}
\resizebox{\columnwidth}{!}{%
\begin{tabular}{lcccc}
\toprule
Method & Class-IL$\uparrow$ & Sel.\ Jaccard$\uparrow$ & Leakage$\downarrow$ & $\Delta U\uparrow$ \\
\midrule
Sequential fine-tuning & $0.4774$ & $0.0023$ & $0.4408$ & $-220$ \\
LoRA merge only & $0.4466$ & $0.0010$ & $0.4756$ & $-301$ \\
$+$ replay & $0.7778$ & $0.0864$ & $0.104$ & $-21.2$ \\
$+$ distillation & $0.7972$ & $0.1752$ & $0.117$ & $-31.9$ \\
PathSelect (full) & $0.8239$ & $0.1613$ & $0.1005$ & $-16.8$ \\
\bottomrule
\end{tabular}}
\end{table}

Table~\ref{tab:selmetrics} shows what accuracy cannot: unprotected selectors retain almost none of their original evidence (Jaccard $0.0023$/$0.0010$) and lose old-task utility by an order of magnitude more than protected ones ($\Delta U$ $-220$/$-301$ vs.\ $-16.8$).
LoRA merging alone shows no systematic difference from full fine-tuning (3/5): merging is the deployment substrate, not the preservation mechanism.
Distillation without the hinge retains slightly more identity ($0.1752$ vs.\ $0.1613$) yet is less accurate and forgets more: preserving \emph{where} the policy looks does not secure \emph{how useful} what it finds is---the gap the utility hinge closes.
Three further interventions raised Jaccard without improving diagnosis: warm-starting task~1 with full-parameter training, damping each fold-in to half, and composing independently trained residuals post hoc.

\textbf{Update-protocol probes.}
Each was decided by a rule fixed before the run.
Warm start and damped merging showed no systematic benefit.
A single adapter never reset lost $3.54$ class-IL points on the flat substrate (4/5) and became a trade-off under the hierarchy (systematic task-IL edge for fresh adapters, systematically lower leakage for the single adapter); fresh adapters are retained.
Post-hoc composition refuted our registered expectation: independently trained residuals, summed once, beat the bare sequential substrate on every metric in every seed, consistent with strong interference in unprotected sequential updates (cf.\ task arithmetic \cite{ilharco2023task})---yet trailed the full method by $13.1$ points (5/5), because independence excludes replay and distillation targets require an accumulated model; sequential consolidation earns its place only together with preservation.
Porting MergeSlide's orthogonal-projection merging \cite{bui2026mergeslide} to our selector did not close this gap ($0.709{\pm}.035$; no systematic difference from plain summation, 3/5; the released implementation's diagonal mask is a no-op and the two variants differ within noise).
With 128 memory entries PathSelect still outperformed replay with 512 (task-IL $+3.06$, 5/5), reported for task-IL and this range only.
A task-conditioned query, a stateful multi-round selector, and a group-level semantic prior failed their registered gates and were removed.



\FloatBarrier
% -----------------------------------------------------------------------------
\section{Discussion and Limitations}
\label{sec:discussion}

PathSelect suggests that continual learning in budgeted perception should distinguish two failure modes that accuracy alone conflates.
A policy may change \emph{where} it looks, and the new evidence may or may not remain \emph{functionally equivalent}.
This is why Jaccard, KD, utility, and the utility hinge should not be collapsed into one notion.
KD preserves a soft behavior distribution; Jaccard measures hard identity after the fact.
Candidate utility $u_i$ teaches which patch is useful now; set utility $U$ evaluates whether the selected evidence remains useful across time.
The utility-preservation hinge uses $U$ as a lower bound, deliberately allowing better alternative evidence.

The study also reveals what is \emph{not} sufficient.
LoRA merging alone does not protect the selector.
Warm start, damped merging, and post-hoc composition each increase selection Jaccard without improving diagnosis (Sec.~\ref{ssec:ablation}), so ``choose the same patches'' is not a complete CL objective: selection stability is not sufficient for diagnostic preservation, and identity and function must be measured separately.
This motivates behavioral and functional preservation as complementary constraints rather than redundant regularizers.

Our claims remain bounded.
The benchmark uses four organ-separated TCGA tasks, one frozen foundation model (CONCH), one main budget ($B=8$), known task boundaries during training, and five seeds.
The tasks are highly separable (a linear probe separates them at 98.21\%), and explicit task conditioning failed its registered gate, so harder within-organ or mixed-task streams may yield different conclusions.
PathSelect is also a one-shot evidence-acquisition policy, not a multi-step navigation agent.
Its task-free shared selector can serve as a foundation for continual evidence-navigation systems, but sequential stateful acquisition is left to future work. The memory is compact but is not claimed to provide a formal privacy guarantee. The hierarchical substrate's advantage and the accuracy increment of distillation and the utility hinge over replay are order-dependent (Table~\ref{tab:chain}); their forgetting reduction holds in 8/10 folds in both orders. All pre-registered comparisons were made on the canonical reverse order; forward-order results are reported as measured. Our reproduction of QPMIL-VL shows twice the fold-to-fold spread of the published value, so single-split rankings on this benchmark should be read with that variance in mind.

% -----------------------------------------------------------------------------
\section{Conclusion}

PathSelect reframes continual WSI learning around a simple but under-measured question: after learning new diagnostic tasks, does the model still know \emph{where to look}?
By freezing the diagnostic pathway, we isolate the evidence selector as the study object.
Within tasks, semantic and counterfactual signals teach the selector to find useful evidence.
Across tasks, a single merged selector is preserved at two complementary levels: distillation maintains selection behavior, while a utility-preservation hinge allows evidence identity to change without allowing usefulness to regress; replay maintains old-task competence.
The resulting method substantially reduces forgetting at zero inference-time memory or routing cost.
More broadly, the results argue that continual learning for budgeted perception should preserve not only what a model predicts, but also the policy by which it acquires the evidence used to predict.

\bibliographystyle{IEEEbib}
\bibliography{refs}

\end{document}
